%------------------------------------------------
\section{Кольцо в форме 8}
%------------------------------------------------
\begin{frame}
	\frametitle{Сохранение поляризации}
	\par Согласно уравнению Т-БМТ, поляризация пучка совершает прецессию в поперечном магнитном поле относительно направления импульса с относительной частотой $\nu_s=\gamma G$
	\newline
	\par Для кольца в форме 8 магнитное поле имеет различное направление в двух поворотных арках. Для частиц любой энергии спиновая прецессия в одной арке полностью компенсируется другой -- режим спиновой прозрачности.
	\newline
	\par Отсутствует необходимость пересечения спиновых резонансов при ускорении протонного поляризованного пучка
\end{frame}
%------------------------------------------------
\begin{frame}
	\frametitle{NICA bypass}
	\centering
	\includegraphics[width=0.65\columnwidth]{4_NICA_bypass.png}
\end{frame}
%------------------------------------------------
\begin{frame}
	\frametitle{Кольцо в форме 8}
	\begin{minipage}[t]{0.4\textwidth}
		\begin{figure}
			\centering
			\includegraphics[width=0.5\linewidth]{"../../6. RuPAC'25 СПб/img/floor_plan_8_shape"}
			\label{fig:floorplan8shape}
		\end{figure}
	\end{minipage}
	\hfill
	\begin{minipage}[t]{0.59\textwidth}
		\vspace{0.4cm}
		\par $L_{mag}=\ 242.8$ см ($1.8$ Тл). Количество магнитов $48$.
		\newline
		\par Контроль $\alpha$ определяется отношением $S/\nu_x$. 
		\newline 
		\par Оптимальным является выбор $S=4$ и равного $\nu_{x,y}=3$.
		\newline
		\par Условия ахромата первого и второго порядка.
	\end{minipage}
	\begin{figure}
		\centering
		\includegraphics[width=0.45\linewidth]{"../../6. RuPAC'25 СПб/img/arc"}
		\label{fig:arc}
	\end{figure}
\end{frame}
%------------------------------------------------
\begin{frame}
	\frametitle{Твисс-функции резонансной структуры кольца в форме 8}
	\begin{figure}
		\centering
		\includegraphics[width=0.9\linewidth]{"../../6. RuPAC'25 СПб/img/8shape Twiss resonant"}
		\label{fig:8shape-twiss-resonant}
	\end{figure}
\end{frame}
%------------------------------------------------
\begin{frame}
	\frametitle{Твисс-функции регулярной структуры кольца в форме 8}
	\begin{figure}
		\centering
		\includegraphics[width=0.9\linewidth]{"../../6. RuPAC'25 СПб/img/8shape Twiss gold"}
		\label{fig:8shape-twiss-gold}
	\end{figure}
\end{frame}
